\paragraph{Primera Ronda}

Se distribuye la pregunta inicial a los expertos, quienes responden de forma independiente.

\begin{center}
\captionof{table}{Ronda 1 - Proyección de clientes por perfil de experto}
\begin{tabular}{|l|c|c|c|}
\hline
\textbf{Ronda 1 - Clientes} & \textbf{Año 1} & \textbf{Año 2} & \textbf{Año 3} \\
\hline
CEO Startup & 8-15 & 23-45 & 50-70 \\
\hline
Gerente de producto & 8-15 & 15-30 & 22-50 \\
\hline
Analista de mercado & 10-20 & 18-44 & 27-75 \\
\hline
Consultor en tecnología & 8-12 & 20-30 & 40-60 \\
\hline
Investigador académico & 8-12 & 16-30 & 30-60 \\
\hline
\end{tabular}
\end{center}

\paragraph{Análisis y Retroalimentación}

Se promedian estos resultados y se vuelven a correr las consultas para dar lugar a las predicciones de la segunda ronda.

\begin{center}
\captionof{table}{Ronda 2 - Proyección de clientes por perfil de experto}
\begin{tabular}{|l|c|c|c|}
\hline
\textbf{Ronda 2 - Clientes} & \textbf{Año 1} & \textbf{Año 2} & \textbf{Año 3} \\
\hline
CEO Startup & 10-15 & 20-30 & 35-50 \\
\hline
Gerente de producto & 10-18 & 18-40 & 30-60 \\
\hline
Analista de mercado & 10-15 & 20-33 & 30-60 \\
\hline
Consultor en tecnología & 12-15 & 25-35 & 21-35 \\
\hline
Investigador académico & 10-15 & 20-35 & 40-60 \\
\hline
\end{tabular}
\end{center}

\paragraph{Consolidación de resultados y escenarios}

Aplicando medianas y rangos intercuartílicos (IQR), se establecen tres escenarios:

\begin{center}
\captionof{table}{Escenarios de negocio consolidados}
\begin{tabular}{|l|c|c|c|c|}
\hline
\textbf{Variable} & \textbf{Pesimista (P25)} & \textbf{Base (Mediana)} & \textbf{Optimista (P75)} & \textbf{Unidad} \\
\hline
Clientes Año 1 & 8 & 12 & 16 & clientes \\
\hline
Clientes Año 2 & 18 & 25 & 35 & clientes \\
\hline
Clientes Año 3 & 30 & 45 & 65 & clientes \\
\hline
Precio anual (ACV) & 10,000 & 12,000 & 15,000 & USD \\
\hline
CAC & 3,500 & 2,800 & 2,200 & USD \\
\hline
Churn anual & 15\% & 10\% & 7\% & \% \\
\hline
Tiempo implementación & 4 & 3 & 2 & meses \\
\hline
\end{tabular}
\end{center}

\paragraph{Análisis de sensibilidad}

Se evalúa el impacto de variaciones en las variables clave sobre los ingresos del Año 3:

\begin{center}
\captionof{table}{Análisis de sensibilidad - Impacto en ingresos Año 3}
\begin{tabular}{|l|c|c|c|}
\hline
\textbf{Variable} & \textbf{Variación} & \textbf{Impacto en ingresos} & \textbf{Sensibilidad} \\
\hline
Precio (ACV) & +/- 20\% & +/- \$108k & Alta \\
\hline
Adquisición clientes & +/- 20\% & +/- \$108k & Alta \\
\hline
CAC & +/- 20\% & +/- \$25k & Media \\
\hline
Churn & +/- 5pp & +/- \$65k & Media-Alta \\
\hline
\end{tabular}
\end{center}

Las variables de mayor impacto son el precio y la capacidad de adquisición de clientes, seguidas por el churn. El CAC tiene menor sensibilidad relativa en los ingresos totales.

\paragraph{Validación contra benchmarks}

Los resultados se validaron contra benchmarks publicados:
\begin{itemize}
    \item \textbf{CAC/ACV ratio:} 0.19-0.23 (benchmark SaaS: 0.15-0.30) - Dentro del rango
    \item \textbf{Churn anual:} 7-15\% (benchmark B2B: 5-20\%) - Dentro del rango
    \item \textbf{Tiempo implementación:} 2-4 meses (benchmark enterprise: 3-6 meses) - Dentro del rango
\end{itemize}

\subsubsection{Proyección Monetaria por Escenarios}

Utilizando los resultados del Delphi simulado, se proyectan tres escenarios financieros. El escenario base utiliza la mediana de las estimaciones (ACV: USD 12,000; implementación: USD 10,000).

\paragraph{Proyección de ingresos por escenarios}

\begin{center}
\captionof{table}{Proyección de ingresos por escenario - Primeros 3 años}
\begin{tabular}{|c|c|c|c|c|}
\hline
\textbf{Año} & \textbf{Escenario} & \textbf{Clientes activos} & \textbf{Nuevos clientes} & \textbf{Ingresos [kUSD]} \\
\hline
1 & Pesimista & 8 & 8 & 160 \\
\hline
1 & Base & 12 & 12 & 264 \\
\hline
1 & Optimista & 16 & 16 & 400 \\
\hline
2 & Pesimista & 18 & 10 & 316 \\
\hline
2 & Base & 25 & 13 & 430 \\
\hline
2 & Optimista & 35 & 19 & 665 \\
\hline
3 & Pesimista & 30 & 12 & 480 \\
\hline
3 & Base & 45 & 20 & 740 \\
\hline
3 & Optimista & 65 & 30 & 1,175 \\
\hline
\end{tabular}
\end{center}

\paragraph{Análisis de punto de equilibrio}

En el escenario base, considerando costos operativos de USD 180k anuales:
\begin{itemize}
    \item \textbf{Punto de equilibrio:} 15 clientes activos (mes 18)
    \item \textbf{Margen bruto:} 85\% (excluyendo costos de implementación)
    \item \textbf{Payback de CAC:} 2.8 meses promedio
\end{itemize}
